\renewcommand{\thechapterdisplay}{Terceiro dia}
\renewcommand{\thechapterrunner}{Existência do Purgatório I}
\chapter[Terceiro dia --- Existência do Purgatório I]{Existência do Purgatório I}
\section*{A Palavra de Deus}

A existência do Purgatório não é somente uma crença piedosa, que somos livres para aceitar ou rejeitar, é um dogma formal, ensinado pela fé, e que devemos professar sob pena de excomunhão. Sim, \textit{é um pensamento santo e salutar}, diz o Antigo Testamento, \textit{o de rezar pelos mortos, afim de que sejam livrados de seus pecados}. Os judeus estavam tão convencidos desta verdade, que eles tinham em seu Ritual uma oração especial que o chefe da família deveria fazer para a libertação dos falecidos antes de se sentarem à mesa.

Mas escutemos o próprio Jesus Cristo: \textit{"Acertai as contas com vosso adversário, enquanto estais em vida: pois, do contrário, vosso adversário o entregará nas mãos do juiz, e o juiz o entregará ao seu ministro que o lançará na prisão, de onde só saireis quando tiverdes pago vossa dívida até o último centavo."} Ora, esse adversário, nos diz Santo Agostinho, é o próprio Deus, o inimigo irreconciliável do pecado. Esse juiz inexorável é Jesus Cristo, que se chama, na Escritura, o juiz dos vivos e dos mortos. Enfim, essa prisão temível é o Purgatório, de onde não se pode sair senão depois de ter satisfeito inteiramente à justiça divina.

Jesus não se contentou em gravar em nossos corações a lembrança do Purgatório, ele nos deu o exemplo; pois, após sua morte, ele desceu ao limbo, para consolar e libertar as almas dos Patriarcas. 

Meu Deus! Eu creio no Purgatório, como em todos os dogmas que Vós ensinais, porque Vós não podeis nem Vos enganar nem me enganar. Eu adoro a equidade de Vossos julgamentos, mesmo nos rigores de Vossa justiça!


\section*{O ensinamento da Igreja}

A fé da igreja não é menos explícita. Veja como a formula o concílio de Trento: \textit{Que seja anátema quem afirma que, após haver recebido a graça da justificação, todo pecador obtém de tal modo a remissão de sua falta e o perdão da pena eterna, que não lhe resta nenhuma dívida temporal a pagar, nem neste mundo nem no outro -- no Purgatório -- antes que lhe seja aberta a entrada do reino dos Céus.} Todos os doutores, gregos e latinos, todos os povos antigos e modernos
professaram a mesma crença.

De acordo com este ponto de fé, a Igreja, mãe terna e compassiva, reza todos os dias pelas almas do Purgatório, e termina cada um de seus ofícios com um clamor de dor e de esperança: \textit{Senhor, dai-lhes o repouso eterno!} Ela obriga todos os seus sacerdotes a pensar nelas, no santo sacrifício, e recomenda a seus filhos oferecer frequentemente a Deus seus votos, suas esmolas e todas as obras satisfatórias pela libertação de seus irmãos falecidos. Enfim, Ela determinou um solene aniversário, onde chama a cristandade inteira em socorro dos fiéis falecidos. É a comovente festa de dois de novembro.

Cristãos, como é consolador pensar que após nossa morte, a Igreja estará por nós em orações, que ela convidará todos os seus fiéis a pedir a Deus nossa libertação, que ela não cessará de rezar e de chorar sobre nosso túmulo, senão quando nos tiver introduzido no seio da Igreja triunfante! Ó Igreja católica, como sois boa mãe! Como conheceis a fraqueza de vossos filhos! E como é doce dizer com o poeta:

\begin{quote}
\textit{\textit{A esperança de que amigos chorarão nossa sorte} \textit{Encanta o instante supremo e consola a morte!}}
\end{quote}


\section*{Exemplo}

Judas Macabeu, esse homem de fé e de coração, a quem o Senhor havia confiado o cuidado de defender Israel e sua lei, Jerusalém e seu templo, acabava de obter uma grande vitória e de pôr em fuga os inimigos de seu Deus e de sua pátria. O primeiro movimento desse guerreiro, não menos piedoso que corajoso, foi dobrar os joelhos para render graças ao Deus dos exércitos. Depois, levantando-se com os seus, ele viu ao seu redor os corpos de seus companheiros de armas que haviam morrido; e então, imbuído de um santo respeito pelos restos inanimados desses bravos guerreiros, Judas os recolheu com cuidado para depositá-los no sepulcro de seus pais. E enfim, pensando nas almas desses mártires da religião e da pátria, ele mandou fazer uma coleta e enviou a Jerusalém doze mil dracmas de prata, a fim de obter um sacrifício pelos pecados dos mortos. Pois ele pensava com sabedoria e piedade na ressurreição, considerando que aqueles que haviam adormecido na religião tinham reservado uma recompensa preciosa.

Eis o que se passava, há dois mil anos, em um campo de batalha da Palestina. Confirmando todas essas coisas tão sóbrias e tocantes, o Espírito de Deus repete pela boca do historiador sagrado: \textit{É, portanto, um pensamento santo e salutar rezar pelos mortos a fim de que sejam libertados de seus pecados.} Assim, Judas Macabeu acreditava no Purgatório e o Espírito Santo louva e confirma essa crença. Infelizmente, onde estão nos dias de hoje os generais do exército capazes de imitar essa piedosa conduta?

\vspace{\baselineskip}\noindent \textbf{Oremos.} Filho submisso de Vossa Igreja, eu creio firmemente, ó meu Deus, na existência do Purgatório. Eu creio porque Vosso Espírito de verdade a revelou, porque Vossos Santos e Vossos Doutores o ensinam. Aumentai minha fé a fim de dilatar minha caridade para com as Almas cativas. Ó Jesus, sede-lhes propício! Senhor, chamai Vossos filhos, nossos irmãos, ao repouso eterno. Que a luz que nunca se apaga brilhe sobre eles! Que descansem em paz!



\begin{center}
\textit{Requiescant in pace!}
\end{center}

Rezai agora uma dezena do terço e as orações no final deste livro.
